% Instant Pot Chicken, Rice & Lentils
% Save as: InstantPot/Chicken_Rice_Lentils.tex
% Add to main file with: % Instant Pot Chicken, Rice & Lentils
% Save as: InstantPot/Chicken_Rice_Lentils.tex
% Add to main file with: % Instant Pot Chicken, Rice & Lentils
% Save as: InstantPot/Chicken_Rice_Lentils.tex
% Add to main file with: % Instant Pot Chicken, Rice & Lentils
% Save as: InstantPot/Chicken_Rice_Lentils.tex
% Add to main file with: \include{InstantPot/Chicken_Rice_Lentils}
% (Uncomment the \chapter line below only if this is the first recipe in a new chapter.)

%\chapter{Instant Pot}
%\minitoc

\begin{recipebox}{Instant Pot Chicken, Rice \& Lentils}{3}{\textasciitilde40 min}

A protein-boosted take on Instant Pot chicken and rice. Red lentils cook
alongside the rice in the same pressure cycle, so there are no extra steps ---
just a bit more stock to cover what the lentils soak up.

\begin{ingredients}
  \item 400\,g skinless chicken breast, diced into bite-sized pieces
  \item 1 cup long-grain white rice, rinsed
  \item \nicefrac{1}{2} cup dried red lentils, rinsed
  \item 2\nicefrac{1}{2} cups chicken stock
  \item 1 medium onion, diced
  \item 3 garlic cloves, minced
  \item 1 tbsp olive oil
  \item 2 tbsp butter
  \item 1 tsp paprika
  \item \nicefrac{1}{2} tsp dried oregano
  \item 1 tsp salt
  \item \nicefrac{1}{4} tsp black pepper
  \item 1 cup frozen peas or mixed vegetables (optional)
  \item 2 tbsp fresh parsley, chopped, to serve (optional)
\end{ingredients}

\begin{method}
  \item \textbf{Saut\'e the aromatics.} Set the pot to Saut\'e. Add the olive
        oil and butter. Once melted, cook the onion for 2--3 minutes until
        softened, then add the garlic and cook for 30 seconds until fragrant.
  \item \textbf{Sear the chicken.} Push the onion to the side and add the
        diced chicken with the paprika, oregano, salt and pepper. Sear for
        1--2 minutes, stirring occasionally --- it only needs colour, not to
        cook through.
  \item \textbf{Add rice, lentils and stock.} Hit Cancel. Sprinkle the rice
        and lentils evenly over the chicken, then pour the stock over the top
        \emph{without stirring}.
  \item \textbf{Pressure cook.} Lid on, valve sealed. Pressure Cook / Manual on
        High for \textbf{10 minutes}.
  \item \textbf{Release.} Allow a 10-minute natural release, then quick-release
        any remaining pressure.
  \item \textbf{Finish.} Stir through the frozen vegetables (if using) --- the
        residual heat warms them through --- and fluff everything together.
  \item Garnish with parsley and serve warm.
\end{method}

\begin{notes}
Stick with red or split lentils --- brown or green lentils need longer and
won't match the rice timing. The extra stock allows for the lentils soaking up
liquid. Dicing the chicken means it cooks through fully under pressure. If you
get a burn warning, scrape any stuck bits off the base before adding the rice
and lentils. For a dairy-free batch, use extra olive oil in place of the butter.
Leftovers reheat and freeze well.
\end{notes}

\end{recipebox}

% (Uncomment the \chapter line below only if this is the first recipe in a new chapter.)

%\chapter{Instant Pot}
%\minitoc

\begin{recipebox}{Instant Pot Chicken, Rice \& Lentils}{3}{\textasciitilde40 min}

A protein-boosted take on Instant Pot chicken and rice. Red lentils cook
alongside the rice in the same pressure cycle, so there are no extra steps ---
just a bit more stock to cover what the lentils soak up.

\begin{ingredients}
  \item 400\,g skinless chicken breast, diced into bite-sized pieces
  \item 1 cup long-grain white rice, rinsed
  \item \nicefrac{1}{2} cup dried red lentils, rinsed
  \item 2\nicefrac{1}{2} cups chicken stock
  \item 1 medium onion, diced
  \item 3 garlic cloves, minced
  \item 1 tbsp olive oil
  \item 2 tbsp butter
  \item 1 tsp paprika
  \item \nicefrac{1}{2} tsp dried oregano
  \item 1 tsp salt
  \item \nicefrac{1}{4} tsp black pepper
  \item 1 cup frozen peas or mixed vegetables (optional)
  \item 2 tbsp fresh parsley, chopped, to serve (optional)
\end{ingredients}

\begin{method}
  \item \textbf{Saut\'e the aromatics.} Set the pot to Saut\'e. Add the olive
        oil and butter. Once melted, cook the onion for 2--3 minutes until
        softened, then add the garlic and cook for 30 seconds until fragrant.
  \item \textbf{Sear the chicken.} Push the onion to the side and add the
        diced chicken with the paprika, oregano, salt and pepper. Sear for
        1--2 minutes, stirring occasionally --- it only needs colour, not to
        cook through.
  \item \textbf{Add rice, lentils and stock.} Hit Cancel. Sprinkle the rice
        and lentils evenly over the chicken, then pour the stock over the top
        \emph{without stirring}.
  \item \textbf{Pressure cook.} Lid on, valve sealed. Pressure Cook / Manual on
        High for \textbf{10 minutes}.
  \item \textbf{Release.} Allow a 10-minute natural release, then quick-release
        any remaining pressure.
  \item \textbf{Finish.} Stir through the frozen vegetables (if using) --- the
        residual heat warms them through --- and fluff everything together.
  \item Garnish with parsley and serve warm.
\end{method}

\begin{notes}
Stick with red or split lentils --- brown or green lentils need longer and
won't match the rice timing. The extra stock allows for the lentils soaking up
liquid. Dicing the chicken means it cooks through fully under pressure. If you
get a burn warning, scrape any stuck bits off the base before adding the rice
and lentils. For a dairy-free batch, use extra olive oil in place of the butter.
Leftovers reheat and freeze well.
\end{notes}

\end{recipebox}

% (Uncomment the \chapter line below only if this is the first recipe in a new chapter.)

%\chapter{Instant Pot}
%\minitoc

\begin{recipebox}{Instant Pot Chicken, Rice \& Lentils}{3}{\textasciitilde40 min}

A protein-boosted take on Instant Pot chicken and rice. Red lentils cook
alongside the rice in the same pressure cycle, so there are no extra steps ---
just a bit more stock to cover what the lentils soak up.

\begin{ingredients}
  \item 400\,g skinless chicken breast, diced into bite-sized pieces
  \item 1 cup long-grain white rice, rinsed
  \item \nicefrac{1}{2} cup dried red lentils, rinsed
  \item 2\nicefrac{1}{2} cups chicken stock
  \item 1 medium onion, diced
  \item 3 garlic cloves, minced
  \item 1 tbsp olive oil
  \item 2 tbsp butter
  \item 1 tsp paprika
  \item \nicefrac{1}{2} tsp dried oregano
  \item 1 tsp salt
  \item \nicefrac{1}{4} tsp black pepper
  \item 1 cup frozen peas or mixed vegetables (optional)
  \item 2 tbsp fresh parsley, chopped, to serve (optional)
\end{ingredients}

\begin{method}
  \item \textbf{Saut\'e the aromatics.} Set the pot to Saut\'e. Add the olive
        oil and butter. Once melted, cook the onion for 2--3 minutes until
        softened, then add the garlic and cook for 30 seconds until fragrant.
  \item \textbf{Sear the chicken.} Push the onion to the side and add the
        diced chicken with the paprika, oregano, salt and pepper. Sear for
        1--2 minutes, stirring occasionally --- it only needs colour, not to
        cook through.
  \item \textbf{Add rice, lentils and stock.} Hit Cancel. Sprinkle the rice
        and lentils evenly over the chicken, then pour the stock over the top
        \emph{without stirring}.
  \item \textbf{Pressure cook.} Lid on, valve sealed. Pressure Cook / Manual on
        High for \textbf{10 minutes}.
  \item \textbf{Release.} Allow a 10-minute natural release, then quick-release
        any remaining pressure.
  \item \textbf{Finish.} Stir through the frozen vegetables (if using) --- the
        residual heat warms them through --- and fluff everything together.
  \item Garnish with parsley and serve warm.
\end{method}

\begin{notes}
Stick with red or split lentils --- brown or green lentils need longer and
won't match the rice timing. The extra stock allows for the lentils soaking up
liquid. Dicing the chicken means it cooks through fully under pressure. If you
get a burn warning, scrape any stuck bits off the base before adding the rice
and lentils. For a dairy-free batch, use extra olive oil in place of the butter.
Leftovers reheat and freeze well.
\end{notes}

\end{recipebox}

% (Uncomment the \chapter line below only if this is the first recipe in a new chapter.)

%\chapter{Instant Pot}
%\minitoc

\begin{recipebox}{Instant Pot Chicken, Rice \& Lentils}{3}{\textasciitilde40 min}

A protein-boosted take on Instant Pot chicken and rice. Red lentils cook
alongside the rice in the same pressure cycle, so there are no extra steps ---
just a bit more stock to cover what the lentils soak up.

\begin{ingredients}
  \item 400\,g skinless chicken breast, diced into bite-sized pieces
  \item 1 cup long-grain white rice, rinsed
  \item \nicefrac{1}{2} cup dried red lentils, rinsed
  \item 2\nicefrac{1}{2} cups chicken stock
  \item 1 medium onion, diced
  \item 3 garlic cloves, minced
  \item 1 tbsp olive oil
  \item 2 tbsp butter
  \item 1 tsp paprika
  \item \nicefrac{1}{2} tsp dried oregano
  \item 1 tsp salt
  \item \nicefrac{1}{4} tsp black pepper
  \item 1 cup frozen peas or mixed vegetables (optional)
  \item 2 tbsp fresh parsley, chopped, to serve (optional)
\end{ingredients}

\begin{method}
  \item \textbf{Saut\'e the aromatics.} Set the pot to Saut\'e. Add the olive
        oil and butter. Once melted, cook the onion for 2--3 minutes until
        softened, then add the garlic and cook for 30 seconds until fragrant.
  \item \textbf{Sear the chicken.} Push the onion to the side and add the
        diced chicken with the paprika, oregano, salt and pepper. Sear for
        1--2 minutes, stirring occasionally --- it only needs colour, not to
        cook through.
  \item \textbf{Add rice, lentils and stock.} Hit Cancel. Sprinkle the rice
        and lentils evenly over the chicken, then pour the stock over the top
        \emph{without stirring}.
  \item \textbf{Pressure cook.} Lid on, valve sealed. Pressure Cook / Manual on
        High for \textbf{10 minutes}.
  \item \textbf{Release.} Allow a 10-minute natural release, then quick-release
        any remaining pressure.
  \item \textbf{Finish.} Stir through the frozen vegetables (if using) --- the
        residual heat warms them through --- and fluff everything together.
  \item Garnish with parsley and serve warm.
\end{method}

\begin{notes}
Stick with red or split lentils --- brown or green lentils need longer and
won't match the rice timing. The extra stock allows for the lentils soaking up
liquid. Dicing the chicken means it cooks through fully under pressure. If you
get a burn warning, scrape any stuck bits off the base before adding the rice
and lentils. For a dairy-free batch, use extra olive oil in place of the butter.
Leftovers reheat and freeze well.
\end{notes}

\end{recipebox}
